\addcontentsline{toc}{section}{СПИСОК ИСПОЛЬЗОВАННЫХ ИСТОЧНИКОВ}

\begin{thebibliography}{99}
	
	\bibitem{harrison_csharp} Ферроне, Х. Изучаем C\# через разработку игр на Unity. 5-е издание / Х. Ферроне. – Санкт-Петербург~: Питер, 2026. – 400 с. – ISBN 978-5-4461-2932-4. – Текст~: непосредственный.
	
	\bibitem{gorelov_winforms} Горелов, С. В. Современные технологии программирования: разработка Windows-приложений на языке C\# / С. В. Горелов ; под ред. П. Б. Лукьянова. – Москва~: Прометей, 2019. – 362 с. – Текст~: непосредственный.
	
	\bibitem{patrashov_math} Патрашов, А. Математическое руководство по созданию компьютерных игр. Справочник / А. Патрашов. – Москва~: Издательские решения, 2016. – 365 с. – ISBN 978-5-4483-2283-9. – Текст~: непосредственный.
	
	\bibitem{wolff_glsl} Вольф, Д. OpenGL 4. Язык шейдеров. Книга рецепсов / Д. Вольф ; пер. с англ. А. Н. Киселёва. – 2-е изд., эл. – Москва~: ДМК Пресс, 2023. – 370 с. – ISBN 978-5-89818-587-9. – Текст~: непосредственный.
	
	\bibitem{snetkov_csharp} Снетков, В. М. Разработка приложений на C\# в среде Visual Studio 2005 / В. М. Снетков. – Москва : Национальный Открытый Университет ИНТУИТ, 2024. – 956 с. – Текст : электронный.
	
	\bibitem{opengl_redbook} Шрайнер, Д. OpenGL. Официальное руководство программиста. 9-е издание / Д. Шрайнер, Г. Селле. – Москва~: ДМК Пресс, 2022. – 1100 с. – ISBN 978-5-97060-987-5. – Текст~: непосредственный.
	
	\bibitem{lengyel_math} Ленгель, Э. Математика для трёхмерной компьютерной графики и разработки игр / Э. Ленгель. – Москва~: Вильямс, 2019. – 524 с. – ISBN 978-5-9909446-2-6. – Текст~: непосредственный.
	
	\bibitem{shreiner_glsl} Шрайнер, Д. OpenGL. Язык шейдеров GLSL. 3-е издание / Д. Шрайнер, Г. Селле. – Москва~: ДМК Пресс, 2021. – 480 с. – ISBN 978-5-97060-945-5. – Текст~: непосредственный.
	
	\bibitem{gregory_engine} Грегори, Дж. Архитектура игрового движка. Введение в разработку игрового программного обеспечения / Дж. Грегори. – Москва~: ДМК Пресс, 2020. – 1200 с. – ISBN 978-5-97060-926-4. – Текст~: непосредственный.
	
	\bibitem{dunn_math} Данн, Ф. Трёхмерная математика для компьютерной графики и разработки игр / Ф. Данн, И. Парберри. – Москва~: ДМК Пресс, 2019. – 560 с. – ISBN 978-5-97060-754-3. – Текст~: непосредственный.
	
	\bibitem{tyukachev_graphics} Тюкачев, Н. А. C\#. Программирование 2D и 3D векторной графики / Н. А. Тюкачев, В. Г. Хлебостроев. – 5-е изд., стер. – Санкт-Петербург~: Лань, 2024. – 320 с. – ISBN 978-5-507-47565-0. – Текст~: непосредственный.
	
	\bibitem{shikin_polygonal} Шикин, Е. В. Компьютерная графика. Полигональные модели / Е. В. Шикин, А. В. Боресков. – Москва~: ДИАЛОГ-МИФИ, 2000. – 461 с. – ISBN 5-86404-139-4. – Текст~: непосредственный.
	
	\bibitem{luchshev_graphics} Luchshev, P. A. Computer graphics: tutorial / P. A. Luchshev, N. G. Gulub. – Kharkiv : KhAI, 2024. – 80 p : официальный сайт – URL: http://dspace.library.khai.edu/xmlui/handle/123456789/8641 (дата обращения: 17.05.2026)
	
	\bibitem{opentk_github} The Open Toolkit Library (OpenTK) : официальный сайт. – URL: https://github.com/opentk/opentk (дата обращения: 19.04.2026).
	
	\bibitem{glb_spec} GLB File Format Specification : официальный сайт. – URL: https://docs.fileformat.com/ru/3d/glb (дата обращения: 01.05.2026).
	
	\bibitem{stl_format} STL Format Description : официальный сайт. – URL: https://industry3d.ru/formaty-dannykh (дата обращения: 11.05.2026).
	
	\bibitem{opengl_spec} OpenGL – The Industry’s Foundation for High Performance Graphics / The Khronos Group : официальный сайт. – URL: https://www.khronos.org/opengl (дата обращения: 01.05.2026).
\end{thebibliography}
